\section{COMMONSENSE NORM BANK 与 Delphi Architecture}

一个伦理判断模型的最重要“参数”可能不是 hidden size，而是 dataset composition。本章把论文、crowd judgments、rules of thumb 与 UNICORN backbone 拆开，说明 Delphi 的能力和偏差如何沿数据供应链进入系统。

\subsection{Norm Bank 是多个数据制度的汇流}

source slide 列出 Social Chemistry、ETHICS、Moral Stories、Scruples 等数据集，并汇入 COMMONSENSE NORM BANK。右侧写出约 1.7M people's ethical judgments on everyday situations。这个数字表示 annotation volume，不表示 1.7M 个独立文化或均衡群体；不同数据集的收集目的和 label schema 也并不相同。

\lecturefigure{slide-31-commonsense-norm-bank-sources.jpg}{COMMONSENSE NORM BANK 的多数据集来源与 1.7M 判断}{官方录像恢复幻灯片 V031；Jiang et al., 2021}

\begin{knowledgebox}{Dataset Aggregation 不是价值中和}
合并数据能扩大 coverage，却不会自动消除 sampling bias。若多个来源都偏向英语、美国网络用户或某种政治语境，数量相加只会更精确地估计该分布。高质量审计应报告来源、时间、人口属性、冲突 label 与 abstention policy。
\end{knowledgebox}

\subsection{会背原则，不等于会应用原则}

GPT-3 morality slide 对比 declarative 与 imperative。模型可以说“把邻居扔进搅拌机是坏事”，却在改写成命令或复杂语境后给出不稳定响应。蓝色 callout 强调 teaching AI with principles 很重要，同时也有 dual-use 与 robustness 问题。读图时应把 failure 定位到 principle application，而不是简单说模型“不知道这条原则”。

\lecturefigure{slide-32-gpt3-morality-declarative-imperative.jpg}{GPT-3 能复述道德原则，却可能无法稳定应用}{官方录像恢复幻灯片 V032；Stanford CS25 Lecture 16}

\subsection{Rules of Thumb 如何连接 Situation 与 Judgment}

规则页先给 situation，再让 annotator 写 rule of thumb，并标注 legality、social norm、moral foundation 等 attributes。完整 slide 写出约 300,000 rules of thumb、104k real-life situations 与 12 structured attributes。结构化属性帮助模型区分“违法”“不礼貌”“伤害”等不同理由，却也把 annotation ontology 变成能力边界。

\lecturefigure{slide-33-rules-of-thumb-300k.jpg}{Rules of Thumb：把现实情境、判断理由与结构属性连接起来}{官方录像恢复幻灯片 V033；Stanford CS25 Lecture 16}

\subsection{UNICORN Backbone 与多能力组合}

架构 slide 将 norm bank 与 UNICORN commonsense model 相加，得到 Delphi。这里的加号不是数学相加，而是 pretraining / fine-tuning pipeline：UNICORN 提供统一 commonsense QA 表示，norm bank 提供 morally salient supervision。系统需要同时处理 language understanding、commonsense inference 和 norm classification。

\lecturefigure{slide-34-delphi-unicorn-architecture.jpg}{COMMONSENSE NORM BANK + UNICORN 形成 Delphi}{官方录像恢复幻灯片 V034；Jiang et al., 2021}

下一页直接回答“Why Delphi is built on top of Unicorn?”：moral reasoning 需要 language、commonsense 与 norms 同时工作。右侧例子把点汉堡、医学建议等情境放在一起，说明伦理判断不是独立 label head，而是依赖世界事实和行为后果。

\lecturefigure{slide-35-why-unicorn.jpg}{为什么 moral reasoning 需要 language、commonsense 与 norms}{官方录像恢复幻灯片 V035；Stanford CS25 Lecture 16}

\begin{warningbox}{Backbone 不能修复 Label Semantics}
更强 encoder 可以提高分布内 generalization，却不会决定“多数票是否公正”“冲突文化如何表达”或“模型何时应拒答”。这些属于 data governance、task design 与 deployment policy，而不是 representation capacity 自动解决的问题。
\end{warningbox}

\subsection{本章小结}

Delphi 的行为来自三层：多来源 norm data、结构化 rules of thumb 与 commonsense backbone。理解这条链后，系统的高分和偏差都不再神秘。下一章转向公开发布后的 adversarial pressure，检验 dataset 内表现能否承受真实用户。

